\section{Related Work}
\label{sec:related-work}

\subsection{LLM Agents for Instructional Design}

The systems closest to ours apply ISD theory to agent architecture.
ISD-Agent-Bench~\cite{isdagentbench} implements four theory-grounded agents:
ReAct-ADDIE (a ReAct loop over ISD tools~\cite{yao2023react}), ADDIE-Agent (a
linear five-phase state graph~\cite{branch2009addie}), Dick-Carey-Agent (the
systems approach with a formative-evaluation loop~\cite{dick1978systematic}), and
RPISD-Agent (rapid prototyping~\cite{tripp1990rapid}). EduPlanner~\cite{zhang2025eduplanner}
uses adversarial generate--evaluate--optimize collaboration among agents;
LessonPlanner~\cite{lessonplanner2024} and GPT-4 teaching-plan
studies~\cite{hu2024teachingplan} improve single artifacts; LearnLM~\cite{learnlm2024}
tunes the model itself for pedagogy. All of these accept a brief and optimize the
draft produced by a design pass, and all treat constructive alignment as something
to be \emph{encouraged} through prompt design or scored after the fact---none
represents alignment as state that the system checks per objective and repairs.

\subsection{Benchmarks and Evaluation}

OpenLearnLM~\cite{openlearnlm2026}, EduBench~\cite{edubench2025},
MathTutorBench~\cite{mathtutorbench2025}, and TutorBench~\cite{tutorbench2025}
evaluate educational capabilities at the level of individual tasks.
ISD-Agent-Bench~\cite{isdagentbench} goes further by scoring both artifact quality
against an ADDIE rubric and the quality of the agent's trajectory, over synthetic
scenarios spanning a Context~$\times$~ISD matrix, with multi-judge grading to
reduce self-preference bias. We adopt it as our evaluation substrate for both
research questions, and additionally report a deterministic \emph{residual
alignment-defect audit} that measures a property the rubric does not directly
score.

\subsection{Consistency Control in Multi-Artifact Generation}

Theoretically, constructive alignment~\cite{biggs1996enhancing} defines the
consistency standard among objectives, activities, and assessments, and the
systems approach~\cite{dick1978systematic} justifies systematic verification.
Bloom's taxonomy~\cite{anderson2001taxonomy} provides the cognitive-level standard
an assessment must match. Technically, verifier and self-repair mechanisms such as
Self-Refine~\cite{madaan2023selfrefine} and Reflexion~\cite{shinn2023reflexion}
demonstrate the value of checking and correcting model output from outside the
model---value that grows as base models saturate and as smaller models are
deployed~\cite{cho2026harness,belcak2025slm}. In existing ISD systems, however,
alignment is only encouraged through prompts or scored after generation; it is
never represented as first-class state---a data structure whose links the system
can check individually, whose affected regions it can mark on change, and through
which it can propagate repairs under control.

\subsection{Synthesis and Positioning}

\Cref{tab:related} summarizes the gap. Current ISD agents optimize a single pass,
are evaluated on artifact and trajectory quality, and lack an explicit
representation of alignment for control; none is evaluated on the SLM setting that
increasingly governs deployment cost. AlignmentGraph-ISD fills the gap by making
alignment the first-class graph state of the harness, decomposing generation
around it so alignment holds by construction, verifying and repairing pedagogical
constraints as per-agent tools, and evaluating the result on both an LLM and an
SLM.

\begin{table}[t]
  \caption{Positioning against representative ISD agents.
  AbC = alignment represented as \emph{checkable state} (not just prompted);
  Ver = per-objective pedagogical verification;
  Rep = automatic repair of detected violations;
  SLM = evaluated on a small language model.}
  \label{tab:related}
  \small
  \begin{tabular}{@{}lcccc@{}}
    \toprule
    Agent & AbC & Ver & Rep & SLM \\
    \midrule
    Baseline (single prompt)               & \ding{55} & \ding{55} & \ding{55} & \ding{55} \\
    ReAct-ADDIE~\cite{isdagentbench,yao2023react}       & \ding{55} & \ding{55} & \ding{55} & \ding{55} \\
    ADDIE-Agent~\cite{branch2009addie}     & \ding{55} & \ding{55} & \ding{55} & \ding{55} \\
    Dick-Carey-Agent~\cite{dick1978systematic} & \ding{55} & (\ding{51}) & \ding{55} & \ding{55} \\
    RPISD-Agent~\cite{tripp1990rapid}      & \ding{55} & \ding{55} & \ding{55} & \ding{55} \\
    EduPlanner~\cite{zhang2025eduplanner}  & \ding{55} & (\ding{51}) & (\ding{51}) & \ding{55} \\
    \textbf{AlignmentGraph-ISD} (ours)     & \ding{51} & \ding{51} & \ding{51} & \ding{51} \\
    \bottomrule
  \end{tabular}
\end{table}
